% ch04 — The Schwarzschild Solution
\chapter{The Schwarzschild Solution}
\label{ch:schwarzschild}

This chapter derives the unique spherically symmetric vacuum solution and
explores its geometric structure. We pay special attention to the coordinate
systems and characteristic radii that the ray tracer uses to set up initial
conditions and detect termination events.

\section{Derivation and Birkhoff's Theorem}
\label{sec:schw-derivation}

We derive the Schwarzschild metric from spherical symmetry and the vacuum
Einstein equations, and state Birkhoff's theorem guaranteeing uniqueness.

\section{Coordinate Systems and the Coordinate Singularity}
\label{sec:schw-coordinates}
\coderef{Spacetime}{Schwarzschild}

We distinguish the coordinate singularity at $r = 2M$ from the physical
singularity at $r = 0$, and introduce Eddington-Finkelstein and other
horizon-penetrating coordinates.

\section{Kerr-Schild Form}
\label{sec:schw-kerr-schild}

We write the Schwarzschild metric in Kerr-Schild form, which is regular at
the horizon and serves as the foundation for the codebase's default coordinate
choice.

\section{Kruskal-Szekeres and Penrose Diagrams}
\label{sec:schw-kruskal}

We construct the maximal analytic extension and use Penrose diagrams to
visualise the global causal structure.

\section{Event Horizon, Photon Sphere, and ISCO}
\label{sec:schw-orbits}

We derive the three characteristic radii---the event horizon at $r = 2M$, the
photon sphere at $r = 3M$, and the innermost stable circular orbit at
$r = 6M$---that govern the ray tracer's termination conditions and accretion
disc models.

\paragraph{Key References.}
The derivation follows \cite{Carroll2004}; Kerr-Schild coordinates are
discussed in \cite{Visser2007}; the orbital analysis follows \cite{Wald1984}.
